\documentclass[11pt]{article}
\usepackage[T1]{fontenc}
\usepackage{textcomp}
\usepackage[margin=0.75in]{geometry}
\usepackage{amsmath,amssymb,amsthm}
\usepackage{array,booktabs}
\usepackage{hyperref}
\setlength{\parindent}{0pt}
\newtheorem{theorem}{Theorem}
\theoremstyle{definition}
\newtheorem{definition}{Definition}
\title{Flow Proof Artifact: Ring-one-left-identity}
\author{Generated by \texttt{flow doc proof}}
\date{Flow Proof Book}
\begin{document}
\maketitle
One is the left multiplicative identity in a ring.\\[0.5em]
\textbf{Source.} dummit-foote — *Abstract Algebra*, §7.1\\[0.5em]
\bigskip
\noindent{\large \textbf{Derived fact 1.} \textit{1 * r = r} \hfill Ring \cdot multiplication \cdot one is the left multiplicative identity}\par
\smallskip\noindent\textit{Coordinate: Ring \cdot multiplication \cdot one is the left multiplicative identity}\par
\smallskip\noindent\textit{Source: dummit-foote}\par
\smallskip\noindent\textit{Built on: one is the left identity, for multiplication on Ring}\par
\medskip
\noindent\textbf{Goal.} 1 * r = r\par
\noindent$\displaystyle \forall r \in Ring\quad 1 \cdot r = r$\par
\medskip
\noindent
\begin{tabular}{@{} >{\raggedright\arraybackslash}p{0.06\textwidth} >{\raggedright\arraybackslash}p{0.47\textwidth} >{\raggedleft\arraybackslash}p{0.41\textwidth} @{}}
\textbf{\#} & \textbf{Proof} & \textbf{Mathematics} \\
\midrule
\textbf{1.} & We prove that one is the left multiplicative identity for multiplication on Ring. & \\[0.45em]
\textbf{2.} & We invoke the definitional clause governing multiplication on Ring: one is the left identity, for multiplication on Ring (instantiated for r). & \\[0.45em]
\textbf{3.} & From step 2, this implies 1 times r equals r. Hence proven. & $\displaystyle 1 \cdot r = r$ \\[0.45em]
\bottomrule
\end{tabular}
\label{thm:Ring-multiplication-one_is_the_left_multiplicative_identity}
\par\medskip\hrule
\end{document}
